\section{Descripción de Requerimientos}

\subsection{Requerimientos de Cantidad de Áreas y Redes}

La topología de la empresa comprende tres segmentos LAN independientes,
cada uno asociado a un área funcional de la organización. La cantidad de
usuarios finales por área, antes de aplicar el factor de crecimiento, es
de 850 en LAN 1, 600 en LAN 2 y 200 en LAN 3. Aplicando un crecimiento
previsto del 10\%, las redes deben dimensionarse para albergar 935, 660 y
220 hosts respectivamente, sin contar los dispositivos de telecomunicaciones
ni los servidores de servicios, que requieren subredes dedicadas.

Adicionalmente, la topología contempla una subred de interconexión entre
routers para cada enlace WAN entre sedes, así como una subred dedicada para
los servidores de servicios. En total, el diseño debe prever al menos seis
subredes: tres para usuarios finales, una para servidores, y las necesarias
para los enlaces punto a punto entre dispositivos de enrutamiento.

\subsection{Requerimientos de Direccionamiento}

El espacio de direcciones asignado al Grupo 1 es \texttt{172.16.192.0/19}
para IPv4 y el prefijo \texttt{2001:dbe:bef9::/48} para IPv6. A partir de
estos bloques se debe derivar un esquema de direccionamiento jerárquico que
satisfaga las necesidades de cada segmento de red, considerando los hosts de
usuario final, los dispositivos de telecomunicaciones activos en cada área y
los servidores de servicios.

La estrategia de asignación de direcciones debe definirse de forma diferenciada
según el tipo de dispositivo: servidores, dispositivos de telecomunicaciones y
equipos de usuario final, justificando en cada caso el mecanismo seleccionado
en función de los requerimientos operativos y de administración de la red.

\subsection{Requerimientos de Servicios}

La red debe garantizar la disponibilidad de tres servicios de aplicación
accesibles desde todos los segmentos de la topología a través del enrutador.
El primer servicio es un servidor web que expone contenido HTTP a los usuarios
de todas las áreas. El segundo es un servidor de correo electrónico que permite
el envío y la recepción de mensajes entre los usuarios de la organización. El
tercero es un servicio de mensajería instantánea implementado sobre sockets TCP,
con soporte para salas de chat, visualización de usuarios conectados mediante el
comando \texttt{/list} e historial local de mensajes exportable a archivo de texto.

Los tres servicios se alojan en un servidor Linux con dirección estática dentro
de la subred de servidores, y deben ser accesibles tanto por IPv4 como por IPv6,
en coherencia con el requisito de doble stack definido para toda la topología.

\subsection{Requerimientos de Seguridad}

\subsubsection{Seguridad de Acceso}

Se debe prevenir el acceso no autorizado a los dispositivos de telecomunicaciones
mediante la configuración de credenciales de acceso tanto en la consola local como
en las líneas de terminal virtual. Las contraseñas deben almacenarse de forma
cifrada en los archivos de configuración de los equipos, evitando que queden
expuestas en texto plano dentro de los respaldos de configuración.

En los servidores Linux, el acceso al sistema operativo debe restringirse mediante
reglas de firewall que limiten los puertos expuestos a los estrictamente necesarios
para la operación de cada servicio. Cualquier puerto no requerido por los servicios
web, de correo o de chat debe permanecer cerrado por defecto.

\subsubsection{Seguridad para Configuración Remota}

La administración remota de los dispositivos de telecomunicaciones debe realizarse
exclusivamente a través de SSH versión 2, deshabilitando el protocolo Telnet en
todas las líneas VTY. Esto garantiza que las sesiones de configuración remota
viajen cifradas a través de la red, protegiéndolas frente a ataques de
interceptación pasiva dentro de los mismos segmentos LAN.

Para los servidores Linux, la configuración remota se realizará también mediante
SSH, preferiblemente con autenticación por clave pública en lugar de contraseña,
reduciendo la superficie de ataque frente a intentos de fuerza bruta.

\subsection{Requerimientos de Enrutamiento}

La topología comprende múltiples sedes que deben mantenerse en comunicación
permanente, por lo que se requiere implementar un protocolo de enrutamiento
dinámico que propague automáticamente las rutas hacia todos los segmentos de
la red. El protocolo seleccionado debe operar tanto sobre IPv4 como sobre IPv6,
garantizando conectividad de extremo a extremo para ambas familias de direcciones
en toda la topología.

El diseño del enrutamiento debe contemplar la convergencia de la red ante posibles
cambios en la topología, así como la correcta redistribución de las rutas hacia
las subredes de servidores y enlaces punto a punto, de modo que cualquier
dispositivo final pueda alcanzar cualquier servicio independientemente del área
en la que se encuentre.